\chapter{Технологическая часть}
\section{Выбор технологий для программной реализации}

Программная реализация разработана с использованием метода рекурсивного дизайна. Данный подход предполагает поэтапное расширение функциональности, начиная с минимально работоспособного ядра~\cite{fowler_refactoring}. На каждом этапе система остаётся в рабочем состоянии, что позволяет проводить непрерывное тестирование~\cite{booch_uml}.

В качестве языка программирования выбран C++ стандарта C++20~\cite{cpp20_features}. Стандарт C++20 предоставляет современные возможности языка, включая концепты и статический полиморфизм для повышения скорости работы. Статическая типизация и прямое управление памятью обеспечивают оптимальное использование вычислительных ресурсов.

Для создания графического интерфейса использован фреймворк Qt версии 6.5~\cite{qt_documentation}. Фреймворк Qt предоставляет инструменты для разработки графических приложений, включая систему виджетов. Архитектура Qt, основанная на модели MVC, обеспечивает разделение данных, представления и логики управления.

Средой разработки выбран Visual Studio Code версии 1.106. Среда обеспечивает эффективную отладку программного кода с поддержкой точек останова и пошагового выполнения~\cite{vscode_debug}.

Для распараллеливания вычислений при растеризации используется библиотека OMP~\cite{openmp60}, при использовании команды \texttt{\#pragma omp parallel for collapse(2)}, где цикл перебирает координаты x и y пространства буфера.

\section{Структура классов}
Ниже представлены UML-диаграммы классов для домена отрисовки и генерации городской среды.
Домены пересекаются только в момент передачи объемных объектов функцией CityMap::exportToScene().
\subsection*{Домен генерации городской среды}

\begin{figure}[H]
	\centering
	\includegraphics[height=1\textwidth,angle=90]{CityClassDiagram}
	\caption{Диаграмма классов домена генерации городской среды}
	\label{fig:CityDomain}
\end{figure}

\subsection*{Структура домена отрисовки}


\begin{figure}[H]
	\centering
	\includegraphics[height=1\textwidth,angle=90]{RendererCore}
	\caption{Диаграмма классов домена отрисовки, ч. 1}
	\label{fig:RendererCore}
\end{figure}
В рамках реализации буфер теней использует ортогональную проекцию, так как в реальном мире расстояние до Солнца предельно велико, по сравнению с масштабами города. Рендер поддерживает наличие более одного источника освещения, в том числе точечных, свет от которых распространяется аналогично перспективной проекции.
\begin{figure}[H]
	\centering
	\includegraphics[height=1\textwidth,angle=90]{Rasterization}
	\caption{Диаграмма классов домена отрисовки, ч. 2}
	\label{fig:Rasterization}
\end{figure}
\begin{figure}[H]
	\centering
	\includegraphics[height=1\textwidth,angle=90]{SceneComposition}
	\caption{Диаграмма классов домена отрисовки, ч. 3}
	\label{fig:SceneComposition}
\end{figure}

\section{Создание городской среды}
На рисунке~\ref{fig:CityCreation} представлена UML-диаграмма последовательности, показывающая взаимодействие классов при генерации городской среды.

\begin{figure}[H]
	\centering
	\includegraphics[width=1\textwidth]{CityCreation}
	\caption{UML-диаграмма последовательности создания городской среды}
	\label{fig:CityCreation}
\end{figure}

\section{Пример работы программы}
Интерфейс программы предполагает настройку направления и оттенка света, возможность задания параметров генерации города, а также настройку качества отображения и угла обзора.
Перемещение камеры достигается за счёт нажатия клавиш W, A, S, D, а также клавиши Shift для ускорения полета. Вращение камерой достигается за счёт передвигания курсора по изображению, при зажатой левой клавише мыши. 

На рисунках~\ref{fig:work_example} и~\ref{fig:work_example2} представлены примеры работы программы. На них представлены различные примеры городов и освещений.
\begin{figure}[H]
	\centering
	\includegraphics[width=1\textwidth]{work_example}
	\caption{Пример работы программы}
	\label{fig:work_example}
\end{figure}
Интерфейс может быть представлен как вкладкой, так и отдельным плавающим окном, как на примере~\ref{fig:work_example2}.
\begin{figure}[H]
	\centering
	\includegraphics[width=1\textwidth]{work_example2}
	\caption{Пример работы программы}
	\label{fig:work_example2}
\end{figure}

\section*{Вывод}
В технологической части определен технологический стек: C++20, Qt 6.5 и Visual Studio Code. Представлены UML-диаграммы классов для доменов генерации городской среды и отрисовки, демонстрирующие модульную архитектуру приложения. Описана реализация буфера теней с ортогональной проекцией для солнечного освещения и поддержкой множественных источников света. Приведена диаграмма последовательности процесса генерации городской среды и примеры работы программы с различными конфигурациями городов и освещения.